\documentclass[../sablona-main.tex]{subfiles}


\begin{document}
	\section{Řešení za pomocí GAMS}
	S využitím kódu:
	\begin{lstlisting}[caption={Model optimalizace výroby stolů}, label={lst:gams_model}]
		* Optimalizace vyroby dvou typu stolu
		* i = index typu stolu (1, 2)
		Set i /1,2/;
		
		* Parametry pro jednotlive typy stolu
		* potrebny_cas(i)  - cas potrebny na vyrobu 1 kusu typu i (hodiny)
		* prodejni_cena(i) - prodejni cena jednoho stolu typu i
		* pocet_noh(i)     - pocet nohou na jednom stole typu i
		* pocet_skel(i)    - pocet sklenenych desek na jednom stole typu i
		* pocet_drev(i)    - pocet drevenych desek na jednom stole typu i
		Parameters 
		potrebny_cas(i)  /1 0.6, 2 1.5/
		prodejni_cena(i) /1 200, 2 350/
		pocet_noh(i)     /1 5,   2 5/
		pocet_skel(i)    /1 0,   2 1/
		pocet_drev(i)    /1 1,   2 0/;
		
		* Cilova funkce: celkovy zisk z prodeje stolu
		Variable zisk;
		
		* Rozhodovaci promenna: pocet vyrobenych stolu jednotlivych typu
		Positive Variable pocet(i);
		
		* Prehled omezeni modelu
		Equations drevo, sklo, nohy, vyplata, cas;
		
		* Omezeni na maximalni dostupne mnozstvi dreva (50 desek)
		drevo..  sum(i, pocet_drev(i)*pocet(i)) =L= 50;
		
		* Omezeni na maximalni dostupne mnozstvi skla (35 desek)
		sklo..   sum(i, pocet_skel(i)*pocet(i)) =L= 35;
		
		* Omezeni na pocet nohou (300 kusu)
		nohy..   sum(i, pocet_noh(i)*pocet(i)) =L= 300;
		
		* Casove omezeni vyroby (maximalne 63 hodin)
		cas..    sum(i, potrebny_cas(i)*pocet(i)) =L= 63;
		
		* Definice zisku jako souctu trzeb za vyrobene stoly
		vyplata.. sum(i, prodejni_cena(i)*pocet(i)) =E= zisk;
		
		* Definice a reseni modelu linearniho programovani
		model stoly /all/;
		solve stoly maximizing zisk using LP;
		
		* Vypis optimalniho zisku a poctu vyrobenych stolu
		display zisk.L, pocet.L;
	\end{lstlisting} 
	
	Po spuštění nám to dalo výsledek 30 skleněných a 30 dřevěných stolů s celkovou tržbou \$16~500.
\end{document}